% =====================================================================
% Introduction  --  TARGET: ~350 words
% =====================================================================
\section{Introduction}
\label{sec:introduction}

\guide{
Use this section to set up \emph{why} the project matters and to map
out the rest of the report. A 1st-class intro covers four things:
\begin{itemize}
  \item \textbf{Industrial context (1 short paragraph):} cite IFR
        statistics on industrial-robot installations~\cite{ifr2024};
        explain why articulated arms with grippers dominate
        pick-and-place; mention hazardous / clean-room environments
        that motivate wireless teleoperation.
  \item \textbf{Brief restated in your own words (1 paragraph):}
        two-unit architecture, $\geq$3 DOF + gripper, wireless link,
        closed-loop joint feedback, manual + autonomous modes.
        Cite~\cite{ee6003brief}.
  \item \textbf{The 60\,\% cap (1 sentence):} acknowledge that core
        requirements alone cap at 60\,\% and state that the report
        explicitly labels core vs advanced work.
  \item \textbf{Roadmap of remaining sections (1 paragraph):} forward-
        reference Sections~\ref{sec:literature} through
        \ref{sec:conclusion} so the marker can navigate.
\end{itemize}
Avoid: jokes, first person (use ``the group'' or passive voice), or
re-stating the abstract verbatim.
}

\placeholder{Industrial robotics has become a cornerstone of modern
manufacturing\ldots[replace with team's own opening paragraph citing
\cite{ifr2024} and \cite{siciliano2009}\ldots]}

\placeholder{The coursework brief~\cite{ee6003brief} requires
\ldots[restate the brief in 4--5 sentences, mentioning all six
mandatory objectives: two units, wireless, 3+\,DOF, closed-loop
feedback, manual + autonomous modes, and performance evaluation].}

\placeholder{The remainder of this report is organised as follows.
Section~\ref{sec:literature} reviews the relevant literature
\ldots[continue the roadmap, naming each subsequent section].}

In line with the brief, every technical section explicitly identifies
which work satisfies a \emph{core} requirement (capped at 60\,\%) and
which constitutes an \emph{advanced} embedded-systems feature
contributing to marks above that ceiling.
